\section{Rowing Biomechanics and Ergometer Rowing}
\label{sec:background_rowing_biomechanics_ergometer}

Rowing is a cyclical, whole-body movement that is commonly described through phases such as the catch, drive, finish, and recovery. During the drive, the rower pushes with the legs, transfers force through the trunk, and completes the stroke with the arms. During the recovery, the sequence is reversed as the rower moves back toward the catch. This rhythm is central to both ergometer rowing and on-water rowing, but the two are not mechanically identical.
% TODO cite: rowing technique / biomechanics overview, e.g. Soper and Hume.
% TODO cite: indoor rowing guide for basic erg terminology if used.

An ergometer constrains the rowing action through a fixed machine frame, handle path, sliding seat, and foot stretchers. It allows the user to perform a rowing-like movement indoors and to generate measurable work, but it does not reproduce all forces and spatial relations of a moving boat. In on-water rowing, the rower interacts with oars, blades, riggers, boat movement, and water resistance. The hands and oars follow a rowing-specific action that includes blade handling, squaring and feathering, and the relation between the shell and the water.
% TODO cite: kinematic comparison of ergometer and on-water rowing.
% TODO cite: foot-stretcher force profiles ergometer vs on-water rowing if used.

This difference is important for the present thesis because the user physically rows on an ergometer while the virtual environment can depict either an ergometer-centred action or an on-water rowing action. A direct one-to-one mapping from physical movement to virtual boat rowing is therefore not possible in a strict biomechanical sense. The system must choose which aspects of the stroke to preserve and which aspects to transform.

The ergometer-congruent representation preserves the indoor rowing apparatus as the main frame of reference. It can align the virtual handle, seat, and body movement with the user's felt setup. This makes the virtual action less like a full on-water rowing scene, but more directly connected to what the user is physically doing.

The boat-centred representation preserves the cultural and sport-specific image of rowing. It can show the user as sitting in a shell, moving through water, and interacting with oars. This may better match what users imagine when they think of rowing as a sport. However, the visual hand and oar movements may need to diverge from the exact physical handle trajectory of the ergometer.

This creates a practical form of fidelity trade-off. Biomechanical or visuomotor fidelity to the ergometer is not the same as contextual fidelity to rowing on water. A system can be faithful to the user's physical input while looking less like the sport, or faithful to the sport depiction while transforming the user's physical input. This trade-off provides the concrete movement basis for the embodiment comparison in this thesis.
